\section{Related Work}
\label{sec:related}

Game-based and gamified pedagogy is a recognizable thread in Brazilian HCI teaching. Experience reports describe gamified courses \cite{miranda2019gamificacao}, card games for teaching interaction-design basics \cite{darin2019goople}, and recent playful artifacts such as gamified workshops and investigative dossiers \cite{santos2025ihcweek,rodrigues2025dossies}. The teaching of HCI evaluation as a specific topic has also received attention, most notably in the experience report by Bim et al.\ on bridging academia, industry, and government in this part of the syllabus \cite{bim2019pontes}. These works confirm that active and playful classroom artifacts are an ongoing community concern, but most target HCI broadly rather than Nielsen's ten heuristics specifically.

A smaller set of digital serious games has emerged to support HCI teaching. Sales and Silva \cite{sales2020jogosserios} conduct a systematic literature review on the area and find it sparse, motivating new contributions. Guimar\~{a}es et al.\ \cite{guimaraes2022avaliacao} compare three Brazilian HCI teaching games --- UsabilityGame, UsabiliCity, and MACteaching --- with respect to usability characteristics, providing both a benchmark and a vocabulary for positioning new artifacts within the same family. Complementary work by Dutra et al.\ \cite{dutra2021mapeamento} maps evaluation methods used for educational serious games, supporting qualitative-first evaluation strategies such as the one adopted in this paper.

The published artifact most closely related to our proposal is Geremias et al.'s face-to-face board game in which players guess each heuristic from drawn or acted cues \cite{geremias2022desvendando}. The work reported here occupies the same niche --- teaching Nielsen's ten heuristics --- but as a web-based, single-player digital experience. % The defining mechanic also differs: rather than guessing heuristics from cues, the player completes the same concrete task in two paired interfaces, one that violates the heuristic and one that respects it, with narration framing each scenario before and after the interaction. To the best of our knowledge, no published Brazilian experience report has covered a web-based, scenario-driven game targeting Nielsen's ten heuristics.
